\section{半角置換}\label{節:半角置換}

正弦と余弦を足す・掛ける・割ることで作った式の積分を考える.このような式は,2変数の\textbf{有理式}$R(X,Y)$の$X,Y$に$\sin x,\cos x$を代入したものなので,一般的な形は
\[
\int R(\sin x,\cos x)\dd x
\]
で与えられる.これを,\,$t$の有理式の積分に帰着させることを目標とする.有理式なら,割り算や部分分数分解で積分できるからである.

さて,何と置けばよいだろうか.候補を挙げてみよう.
\begin{enumerate}
\item $t=\sin x$と置くと,\,$\cos x=\pm\sqrt{1-t^2}$となり,根号が残る.
\item $t=\cos x$と置いても,同じく\,$\sin x$が根号で残る.
\item $t=\tan x$と置くと,\,$\cos^2x=\frac1{1+t^2}$,\,$\sin x\cos x=\frac t{1+t^2}$は有理式になる.しかし$\sin x$単独は$\pm\frac{t}{\sqrt{1+t^2}}$となり,根号が残る.
\end{enumerate}
3では,\,$\cos^2x$と$\sin x\cos x$は有理式になった.したがって,2倍角の公式から
\[
\sin2x=2\sin x\cos x=\frac{2t}{1+t^2},\qquad
\cos2x=2\cos^2x-1=\frac{1-t^2}{1+t^2}
\]
も$t=\tan x$の有理式である.つまり,正接は2倍の角の正弦・余弦を有理式で表す.そこで角を半分にして,\,$\tan\frac x2$を使えばよい.

$t=\tan\frac x2$と置く.2倍角の公式で分子を書き,分母の$1=\cos^2\frac x2+\sin^2\frac x2$で割ると
\begin{equation}\label{式:半角置換}
\begin{aligned}
\sin x&=\cfrac{2\sin\cfrac x2\cos\cfrac x2}
 {\cos^2\cfrac x2+\sin^2\cfrac x2}
 =\frac{2t}{1+t^2},\\
\cos x&=\cfrac{\cos^2\cfrac x2-\sin^2\cfrac x2}
 {\cos^2\cfrac x2+\sin^2\cfrac x2}
 =\frac{1-t^2}{1+t^2},\\
\frac{\dd t}{\dd x}&=\frac12\sec^2\frac x2
 =\frac{1+t^2}{2},\qquad
\dd x=\frac{2}{1+t^2}\dd t
\end{aligned}
\end{equation}
を得る.最初の2式では,中央の分数の分子と分母を$\cos^2\frac x2$で割った.

ここで1つ断りをいれておこう.この割り算と$\tan\frac x2$の定義には,\,$\cos\frac x2\ne0$,すなわち$x\ne(2k+1)\pi$($k$は整数)が必要である.そこで,この置換は$(2k-1)\pi<x<(2k+1)\pi$の1つの区間の中で使う.この区間では$t$は実数全体を狭義単調増加して動き,逆向きには
\[
x=2\arctan t+2k\pi
\]
と書ける.とくに$-\pi<x<\pi$では$x=2\arctan t$である.

この置換を\textbf{半角置換},または\textbf{ワイエルシュトラス置換}という.正弦・余弦・$\dd x$のすべてが$t$の有理式になるので,置換後は有理式の積分になる.

\begin{bookpoint}
\textbf{POINT}\par
$t=\tan\frac x2$と置き,\,$\sin x=\frac{2t}{1+t^2}$,\,$\cos x=\frac{1-t^2}{1+t^2}$,\,$\dd x=\frac{2}{1+t^2}\dd t$を代入して$t$の有理式にする.
\end{bookpoint}

\begin{example}{半角置換}
$\int_0^{\frac\pi2}\frac1{2+\sin x}\dd x$を求めよ.
\end{example}

\paragraphbf{方針}
分母$2+\sin x$の微分$\cos x$は分子に現れないので,\,$t=\tan\frac x2$と置いて$t$の有理式にする.最後は平方完成して$\arctan$で積分する.

\paragraphbf{解答}
$\frac\pi{3\sqrt3}$

\paragraphbf{解法}
$t=\tan\frac x2$と置くと,\,$x=0,\frac\pi2$は$t=0,1$へ移る.式\eqref{式:半角置換}を代入すると
\begin{align*}
\int_0^{\frac\pi2}\frac1{2+\sin x}\dd x
&=\int_0^1\cfrac1{2+\cfrac{2t}{1+t^2}}\cdot
 \frac2{1+t^2}\dd t\\
&=\int_0^1\frac{2}{2(1+t^2)+2t}\dd t\\
&=\int_0^1\frac1{t^2+t+1}\dd t
\end{align*}
となる.分母を平方完成すると$t^2+t+1=\left(t+\frac12\right)^2+\frac34$なので
\begin{align*}
\int_0^1\frac1{t^2+t+1}\dd t
&=\frac2{\sqrt3}\left[\arctan\frac{2t+1}{\sqrt3}\right]_0^1\\
&=\frac2{\sqrt3}\left(\frac\pi3-\frac\pi6\right)\\
&=\frac\pi{3\sqrt3}
\end{align*}
を得る.

\paragraphbf{区間をまたぐ場合}
もし積分区間が$[0,2\pi]$なら,\,$x=\pi$で$t$が定義されなくなる.端点だけを$t=\tan0=0$,\,$t=\tan\pi=0$へ換えて答えを$0$としてはいけない.元の被積分関数は$\pi$で連続でも,置換に使った$t$はそこで途切れる.この場合には$\pi$の前後に分け,\,$t\to+\infty$と$t\to-\infty$を別々に扱う.章末問題で確かめる.

\paragraphbf{正割の積分の別解}
前節の$\int\sec x\dd x$も,半角置換で解ける.\,$-\pi<x<\pi$,\,$x\ne\pm\frac\pi2$で$t=\tan\frac x2$と置くと,\,$\sec x=\frac{1+t^2}{1-t^2}$なので
\begin{align*}
\int\sec x\dd x
&=\int\frac{1+t^2}{1-t^2}\cdot\frac{2}{1+t^2}\dd t\\
&=\int\left(\frac1{1+t}+\frac1{1-t}\right)\dd t\\
&=\log\left|\frac{1+t}{1-t}\right|+C
\end{align*}
となる.一方,\,$\sec x+\tan x$を$t$で書くと
\begin{align*}
\sec x+\tan x
&=\frac{1+\sin x}{\cos x}\\
&=\cfrac{1+\cfrac{2t}{1+t^2}}{\cfrac{1-t^2}{1+t^2}}\\
&=\frac{(1+t)^2}{(1+t)(1-t)}
=\frac{1+t}{1-t}
\end{align*}
である.前節で結果からお見せした$\sec x+\tan x$は,半角置換で書けば$\frac{1+t}{1-t}$という簡単な有理式なのだ.
